\section{Method}
\label{sec:method}

\noindent\textbf{Problem setup.}
Let $x$ denote a street-view image and $a$ a perceptual attribute such
as safety. Our explanatory atom is a \emph{lever intervention}
$l=(c,s,d,\tau)$, where $c$ is a lever concept, $s$ the grounded scene
support, $d$ the intervention direction, and $\tau$ a constrained edit
template. For image $x$, a planner grounds a bounded intervention
vocabulary to the visible scene, yielding candidates $L(x)$. Each lever
is tested independently against the unedited original; edits are never
composed sequentially.

\noindent\textbf{Phase 1: Lever candidate construction.}
Following the 9-page formulation, the intervention vocabulary is
\emph{editor-feasible}: each concept must satisfy three criteria,
\textbf{(i)} empirical grounding in urban-perception
literature~\cite{jacobs1961,newman1972,li2015greenery,jiang2018brokenwindows,portnov2020lighting},
\textbf{(ii)} semantic distinctness from other concepts, and
\textbf{(iii)} prompt-only editability without masks or geometry
redesign. The current pilot uses twelve concepts in four families:
Physical Maintenance, Environmental Amenity, Visual Legibility, and
Mobility Infrastructure. The full vocabulary is listed in
Appendix~\ref{sec:intervention_vocabulary}.

\noindent\textbf{Phase 2: Bounded stochastic intervention.}
Each candidate is edited independently from the original image with a
bounded retry budget. A candidate enters the retained set
\begin{equation}
  V(x) = \{ l_i \in L(x) \mid \exists\, \tilde{x}_i \text{ passing the validity audit} \}
\end{equation}
only if at least one edited sample passes four checks: same-place
preservation, locality, realism, and plausibility. This validity gate is
the key design choice. Generative editors can introduce non-target
changes, so failed edits are excluded from attribution rather than being
treated as noisy evidence. We report valid-edit coverage as
$|V(x)|/|L(x)|$, but interpret it as a property of edit realizability
under bounded prompt control rather than an urban-perception signal in
itself.

\noindent\textbf{Phase 3: Auxiliary scoring.}
For each retained edit $\tilde{x}_i$, we compute an auxiliary shift
\begin{equation}
  \Delta_i = f_a(\tilde{x}_i) - f_a(x),
\end{equation}
where $f_a$ is a pretrained Place Pulse safety scorer. This is strictly
an auxiliary proxy, not the main estimand, so $\Delta_i$ is used only
as a prioritization signal. The main reported set is
\begin{equation}
  E_{\mathrm{aux}}(x,a) = \{ l_i \in V(x) \mid \Delta_i > \theta_{\mathrm{aux}} \},
\end{equation}
with $\theta_{\mathrm{aux}}=0.1$ in this pilot. The method therefore
separates three questions: which levers can be grounded, which can be
realised cleanly, and which retained edits look promising enough to send
to human pairwise evaluation.%
\footnote{This is constrained interventional search, not causal
identification or pixel-level decomposition. Distinct levers are tested
independently against the same original image and are not assumed to add
linearly.}
